%%%%%%%%%%%%%%%%%%%%%%%%%%%%%%%%%%%%%%%%%%%%%%%%%%%%%%%%%%%%%%%%%%%%%%%%%%%%%%%%%
\documentclass{article}
% CREDITS TO JUSTIN KIM (idk where he got this from)

%%%%%%%%%%%%%%%%%%%%%%%%%%%%%%%%%%%%%%%%%%%%%%%%%%%%%%%%%%%%%%%%%%%%%%%%%%%%%%%%%%%%
%%% Default packages:

\usepackage[margin=1in]{geometry} 
\usepackage{amsmath,amsthm,amssymb,amsfonts, fancyhdr, color, comment, graphicx, environ}
\usepackage{xcolor}
\usepackage{mdframed}
\usepackage{bm}
\usepackage[shortlabels]{enumitem}
\usepackage{mathtools}
\usepackage{listings}
\usepackage{stmaryrd}
\usepackage{indentfirst}
\usepackage{hyperref}


%%%%%%%%%%%%%%%%%%%%%%%%%%%%%%%%%%%%%%%%%
%%% Pacakages I've added:

%% For Math 300 (Winter 2021):

\usepackage{mathdots} % for \iddots

\usepackage[ruled]{algorithm2e} % Algorithms
% NOTE: FIND A BETTER ALGORITHM PACKAGE, OR ATLEAST LEARN HOW TO USE THIS ONE BECAUSE DOUBLE INDENTING IS A FUCKING NIGHTMARE (or just import from mathcha?)
%% Example algorithm:
% \begin{center}
% 	\begin{minipage}{0.5\linewidth} % Adjust the minipage width to accomodate for the length of algorithm lines
% 		\begin{algorithm}[H]
% 			\KwIn{$(a, b)$, two floating-point numbers}  % Algorithm inputs
% 			\KwResult{$(c, d)$, such that $a+b = c + d$} % Algorithm outputs/results
% 			\medskip
% 			\If{$\vert b\vert > \vert a\vert$}{
% 				exchange $a$ and $b$ \;
% 			}
% 			$c \leftarrow a + b$ \;
% 			$z \leftarrow c - a$ \;
% 			$d \leftarrow b - z$ \;
% 			{\bf return} $(c,d)$ \;
% 			\caption{\texttt{FastTwoSum}} % Algorithm name
% 			\label{alg:fastTwoSum}   % optional label to refer to
% 		\end{algorithm}
% 	\end{minipage}
% \end{center}



%%%%%%%%%%%%%%%%%%%%%%%%%%%%%%%%%%%%%%%%%%%%%%%%%%%%%%%%%%%%%%%%%%%%%%%%%%%%%%%%%%%%
%%% Default commands:

\renewcommand{\vec}[1]{\mathbf{#1}}
	
\newcommand{\WidestEntry}{$lon_1$}%
\newcommand{\SetToWidest}[1]{\makebox[\widthof{\WidestEntry}]{#1}}%
\newcommand\tab[1][0.61cm]{\hspace*{#1}}
\newcommand{\nats}{\mathbb{N}}
\newcommand{\rats}{\mathbb{Q}}
\newcommand{\reals}{\mathbb{R}}
\newcommand{\Z}[1]{\mathbb{Z}_{#1}}
\newcommand{\BigO}[1]{\mathcal{O}(#1)}
\newcommand{\seq[1]}{(#1_n)}
\newcommand{\subseq[1]}{(#1_{n_k})}
\newcommand{\Lim}[2]{\lim \limits _{#1 \to #2}}
\newcommand{\Min}[2]{\min \{#1, #2\}}
\newcommand{\inv}{^{-1}}
\newcommand{\h}{^\text{th}}
\newcommand{\lrangle}[1]{\langle #1 \rangle}
\newcommand{\abs}[1]{\left\lvert #1 \right\rvert}

\DeclarePairedDelimiter{\ceil}{\lceil}{\rceil}
\DeclarePairedDelimiter{\floor}{\lfloor}{\rfloor}
\DeclareMathOperator{\supp}{supp}
\DeclareMathOperator{\rad}{rad}
\DeclareMathOperator*{\argmin}{arg\,min}
\DeclareMathOperator*{\argmax}{arg\,max}
\DeclareMathOperator*{\Var}{Var}
\DeclareMathOperator*{\Cov}{Cov}
\DeclareMathOperator*{\Corr}{Corr}
\DeclareMathOperator*{\Aut}{Aut}
\newcommand{\prob}[1]{\section*{Problem #1}}


%%%%%%%%%%%%%%%%%%%%%%%%%%%%%%%%%%%%%%%%%
%%% Commands I've added:

%% For Math 300 (Winter 2021):

\newcommand{\lrbrace}[1]{\{ #1 \}}
\newcommand{\powerset}{\mathcal{P}}
\newcommand{\ints}{\mathbb{Z}}

% Source/inspiration: https://tex.stackexchange.com/a/42728:
\newcommand{\numberthis}{\addtocounter{equation}{1}\tag{\theequation}\label{\theequation}}
    % Within an `align*` environment, put `\numberthis` after a line to number it. 
    % Access it with `\eqref{ [number of equation] }`
\newcommand{\numberthiswith}[1]{\addtocounter{equation}{1}\tag{\theequation}\label{#1}}
    % Within an `align*` environment, put `\numberthiswith{ [your_label] }` after a line to number it. 
    % Access it with `\eqref{ [your_label] }`



%%%%%%%%%%%%%%%%%%%%%%%%%%%%%%%%%%%%%%%%%%%%%%%%%%%%%%%%%%%%%%%%%%%%%%%%%%%%%%%%%%%%
%%% Default formatting:

\hypersetup{
    colorlinks=true,
    linkcolor=blue,
    filecolor=magenta,      
    urlcolor=blue,
}

\setlength{\parindent}{0cm}
\setlength{\parskip}{6pt}

\pagestyle{fancy}

% Modified version of problem environment below
% \newenvironment{problem}[2][Problem]
%     { \begin{mdframed}[backgroundcolor=gray!5] \textbf{#1 #2} \\}
%     {  \end{mdframed}}
% \newenvironment{solution}{\textbf{Solution}\\}


%%%%%%%%%%%%%%%%%%%%%%%%%%%%%%%%%%%%%%%%%
%%% Formatting I've added:

% Grey boxes for problem statements (note that I don't have a "solution" section):

% Problem environment, but shows "([input_char])" instead of "Problem [input_char]"
\newenvironment{problem}[2][]
    { \begin{mdframed}[backgroundcolor=gray!5] \textbf{#1 (#2)}}
    {  \end{mdframed}}
% Problem environment, but no "([input_char])" at all
\newenvironment{problem*}
    { \begin{mdframed}[backgroundcolor=gray!5] \\}
    {  \end{mdframed}}



%%%%%%%%%%%%%%%%%%%%%%%%%%%%%%%%%%%%%%%%%%%%%%%%%%%%%%%%%%%%%%%%%%%%%%%%%%%%%%%%%%%%
% Misc things I've added




%%%%%%%%%%%%%%%%%%%%%%%%%%%%%%%%%%%%%%%%%%%%%
% Fill in the appropriate information below
\lhead{Zain Kamal}
\rhead{Math 300 Winter 2021} 
% \chead{\textbf{Homework 2}} % Moved to document

\usepackage{tikz}
\usepackage{tikz-among-us}

\begin{document}
\chead{\textbf{Homeworm 8}}
%%%%%%%%%%%%%%%%%%%%%%%%%%%%%%%%%%%%%%%%%%%%%%%%%%%%%%%%%%%%%%%%%%%%%%%%%%%%%%%%%

get ready for the \begin{tikzpicture}
    \amongUsI[scale=0.1]{red}{gray}
\end{tikzpicture} iest proofs of your lifetime. i am very sorry
\section{}

\begin{problem*}
There is an error in Lemma 8.7, similar to the last time I pointed out my own error. Find it, fix the statement of the lemma, and amend the proof to adapt to this. Discuss in Discord.
\end{problem*}

If $\{b\}$ were an "infinite sequence" of empty sequences it would just be empty, so we'd run into the problem where the ``then there would be no elements from $b_{N+1}$ onwards'' condition for contradiction isn't satisfied. 

So we must specify at the beginning that ``...let $B$ be the set of \textbf{non-empty} finite sequences on $A$.''

\
\hline
%%%%%%%%%%%%%%%%%%%%%%%%%%%%%%%%%%%%%%%%%%%%%%%%%%%%%%%%%%%%%%%%%%%%%%%%%%%%%%%%%
\section{}

\begin{problem*}
\textbf{(Pigeonhole Principle).}
\end{problem*}

\begin{problem} a
Show that there is no injection from $[n+1]$ to $[n]$, for any nonnegative integer. 
\end{problem}

We know $|[ n+1] | >|[ n] |$, so $f:[ n+1]\rightarrow [ n]$ can't be injective.\qed

\

\begin{problem} b
$\dagger$ Suppose I give you 100 naturals. Prove that it is always possible to pick a subset of size 15 such that the difference of any two is a multiple of 7.
\end{problem}

no
\begin{tikzpicture}
    \amongUsI[scale=0.1]{cyan}{gray}
\end{tikzpicture}


\
\hline
%%%%%%%%%%%%%%%%%%%%%%%%%%%%%%%%%%%%%%%%%%%%%%%%%%%%%%%%%%%%%%%%%%%%%%%%%%%%%%%%%
\section{}

\begin{problem*}
Exhibit a bijection from $\mathbb{Z}$ to $\mathbb{Q}$.
\end{problem*}

From proposition 8.8, we know that $|\mathbb{N} |=|\mathbb{Q}^{+} |$. We can apply the same reasoning from the proposition to negative numbers to find $|-\mathbb{N} |=|\mathbb{Q}^{-} |$. Because all sets so far are countably infinite, we can add a single element (0) to both without changing cardinality*. Therefore $|\mathbb{N} |=|\mathbb{N} \cup \{0\} |=|\mathbb{Q}^{+} |=|\mathbb{Q}^{+} \cup \{0\} |$. Knowing that $\mathbb{N} \cup ( -\mathbb{N}) \cup \{0\} =\mathbb{Z}$ and $\mathbb{Q}^{+} \cup \mathbb{Q}^{-} \cup \{0\} =\mathbb{Q}$, we can finally reason that $|\mathbb{Z} |=|\mathbb{Q} |$, and thus there is a bijection from $\mathbb{Z}$ to $\mathbb{Q}$.

\

(*Not sure if we have to prove this but here:)

\textit{Claim: }If we add/remove a finite number of elements from a countably infinite set, it is still countably infinite. 

\textit{Proof: }Let $\mathcal{A}$ be a countably infinite set and choose some $A\subseteq \mathcal{A}$ where $A$ must be a finite set. Let $B=\mathcal{A}\setminus A$. Suppose for the sake of contradiction that $B$ is a finite set. Then $B\cup A =\mathcal{A}$, which would imply $A$ is also a finite set which is a contradiction. Thus $B$ is a countably infinite set.

\
\hline
%%%%%%%%%%%%%%%%%%%%%%%%%%%%%%%%%%%%%%%%%%%%%%%%%%%%%%%%%%%%%%%%%%%%%%%%%%%%%%%%%
\section{}

\begin{problem*}
Show that if a set $A$ has a bijection to a subset of $\mathbb{N}$, then it is countable. Conclude that if $|A|\le |\mathbb{N}|$, then it is countable. 
\end{problem*}

Let $N\subseteq \mathbb{N}$. A bijection between $A$ and $N$ implies $|A|=|N|$. 
\begin{itemize}
\item \underline{Case 1:} If $N\subset \mathbb{N}$, then $|A|< |\mathbb{N} |$, meaning $A$ is a finite set and therefore countable. 
\item \underline{Case 2:} If $N=\mathbb{N}$, then $|A|=|\mathbb{N} |$, meaning $A$ is equinumerous with the naturals (countably infinite) and therefore countable. 
\end{itemize}

Therefore, $A$ is countable for $|A|\leq |\mathbb{N} |$.


\
\hline
%%%%%%%%%%%%%%%%%%%%%%%%%%%%%%%%%%%%%%%%%%%%%%%%%%%%%%%%%%%%%%%%%%%%%%%%%%%%%%%%%
\section{}

\begin{problem} a
Suppose $A_1,A_2$ are countable sets. Show that there is an injection $f:A_1\times A_2\to \mathbb{N}$. [Hint: Fundamental theorem of Arithmetic.]
\end{problem}

Let $f( x,y) =2^{x} \cdot 3^{y}$. By the Fundamental theorem of Arithmetic, for all $n\in \mathbb{N}$, there exist unique $x_{0} ,y_{0} \in \mathbb{N}$ such that $n=f( x_{0} ,y_{0})$. 

Now choose some $a,c\in A_{1}$ and $b,d\in A_{2}$ such that $f( a,b) =f( c,d)$. Some algebraic manipulation reveals
\begin{align*}
f( a,b) & =f( c,d)\\
2^{a} \cdot 3^{b} & =2^{c} \cdot 3^{d}\\
2^{a-c} \cdot 3^{b-d} & =1,
\end{align*}
which implies $a=c$ and $b=d$. This maps to a unique natural number, so $f$ is injective.

\

\begin{problem} b
Generalize to any countable $A_1\times A_2\times \cdots \times A_n$.
\end{problem}

Let $P$ be a sequence of the first $n$ prime numbers indexed by $[ n]$ ($P_{1} =2,\ P_{2} =3,\dotsc $). Let $\vec{x} \in A_{1} \times A_{2} \times \cdots \times A_{n}$ where $\vec{x}$ is also indexed by $[ n]$ ($x_{1} \in A_{1} ,\dotsc $). By the Fundamental theorem of Arithmetic, for all $n\in \mathbb{N}$, there exist unique $\vec{x}$ such that $n=f(\vec{x}) =\prod _{i=1}^{n}( P_{i})^{x_{i}}$. This is injective by similar reasoning as part (a). 

\

\begin{problem} c
Use parts (a) and/or (b) to get another proof that $\mathbb{Q}$ is countable.
\end{problem}

Allowing negative exponents in the function in part (b) allows $f$ to represent any positive real number. This is because if any natural number $n=f(\vec{x}) =\prod _{i=1}^{n}( P_{i})^{x_{i}}$, then any rational number in $( 0,1]$ can be represented with $\prod _{i=1}^{n}( P_{i})^{-x_{i}} =\prod _{i=1}^{n}\frac{1}{( P_{i})^{x_{i}}}$. We denote the function that can map onto any positive real number $f'$.



Thus $f':\mathbb{N}^{n}\rightarrow \mathbb{Q}^+$. As shown in part (a), for all $n\in \mathbb{N}$, $|\mathbb{N}^{n} |=|\mathbb{N} |$. As shown in question (3), $|Q^{+} |=|Q^{+} \cup Q^{-} \cup \{0\} |=|Q|$. Thus $|\mathbb{N} |=|\mathbb{Q} |$, so $\mathbb{Q}$ is countable.

% Allowing negative exponents in the function in part (b) allows $f$ to represent any positive real number (we denote this new function $f'$). Thus $f':\mathbb{N}^{n}\rightarrow \mathbb{Q} +$. As shown in part (a), for all $n\in \mathbb{N}$, $|\mathbb{N}^{n} |=|\mathbb{N} |$. As shown in question (3), $|Q^{+} |=|Q^{+} \cup Q^{-} \cup \{0\} |=|\mathbb{Q}|$. Thus $|\mathbb{N} |=|Q|$, so $\mathbb{Q}$ is countable.

\
\hline
%%%%%%%%%%%%%%%%%%%%%%%%%%%%%%%%%%%%%%%%%%%%%%%%%%%%%%%%%%%%%%%%%%%%%%%%%%%%%%%%%
\section{}

\begin{problem} a
Suppose that $\mathcal{A}$ is a countable collection of countably infinite sets. That is, there are countably many sets contained in $\mathcal{A}$. Show that $\bigcup_{A\in\mathcal{A}} A$ is countable.
\end{problem}

Denote the countably infinite sets in $\mathcal{A}$ as $\mathcal{A} =( A_{1} ,A_{2} ,A_{3} ,\cdots )$, where the elements of $A_{i}$ are indexed as $( A_{i,1} ,A_{i,1} ,\dotsc )$. Put each element of $\mathcal{A}$ in a row of a $2\times 2$ grid like such:
\begin{equation*}
\begin{matrix}
A_{1,1} & A_{1,2} & A_{1,3} & \cdots \\
A_{2,1} & A_{2,2} & A_{2,3} & \cdots \\
A_{3,1} & A_{3,2} & A_{3,3} & \cdots \\
\vdots  & \vdots  & \vdots  & \ddots 
\end{matrix}
\end{equation*}
We end up with an $\mathbb{N} \times N$ grid, where $N\subseteq \mathbb{N}$. We know that there exists an injection $f$ such that $\mathbb{N} \times \mathbb{N}\rightarrow \mathbb{N}$. Thus, we can conclude $\left| \bigcup\nolimits _{A\in \mathcal{A}} A\right| \leq |\mathbb{N} |$, so the joint union is countable.

\

\begin{problem} b
Deduce that the collection of all finite subsets of $\mathbb{N}$ is countable. 
\end{problem}

We will proceed by induction. Let $S( n)$ be the set of subsets of $\mathbb{N}$ such that each subset has cardinality less than or equal to $n$. We know $|S( 0) |=|\{\varnothing \} |=1$, so it is countable. Similarly, $|S( 1) |=|S( 0) |+|\mathbb{N}^{1} |$, which is countable as proven in question 3.

Suppose $S( k)$ is countable for some $k\in \mathbb{N}$. Then $|S( k+1) |=|S( k) |+|\mathbb{N}^{k+1} |$. Both terms are countably finite (there is an injection $\mathbb{N}^{k+1}\rightarrow \mathbb{N}$ which implies $|\mathbb{N}^{k+1} |\leq |\mathbb{N} |$ as prove in 5b). The sum of countable infinities is a countable infinity (as proven in 3), so $S( k+1)$ is countable. 

\
\hline
%%%%%%%%%%%%%%%%%%%%%%%%%%%%%%%%%%%%%%%%%%%%%%%%%%%%%%%%%%%%%%%%%%%%%%%%%%%%%%%%%
\end{document}